\section{Evidence Ledger and Independent Verification}

The evidence ledger makes claims falsifiable after execution. It does not prevent all compromise; it provides a cryptographic consistency relation between events, plans, receipts, and externally pinned history.

\subsection{Per-subject event chain}

Each event stores sequence, event ID, subject, type, UTC time, actor, optional session/turn/record, canonical payload, previous hash, and event hash. For subject $u$,

\begin{equation}
h_i=H(\mathrm{id}_i,u,\mathrm{type}_i,t_i,a_i,s_i,r_i,\mathrm{payload}_i,h_{i-1}),
\label{eq:eventhash}
\end{equation}

with $h_0=\bot$. Global sequence numbers order database insertion, while chain links are per subject. The verifier recomputes every event preimage and checks continuity.

Engine-generated memory events contain message/query/selector digests and structural metadata rather than raw values. Action events similarly carry plan, result, observation, and error digests while raw payloads live elsewhere. Low-level custom action logging can still accept arbitrary caller payloads, so deployment code must preserve this convention.

\subsection{Checkpoints}

A checkpoint stores, for every subject, the current sequence, head hash, and event count, then hashes the checkpoint document itself. Verification applies a containment rule: the pinned sequence must exist and have the pinned hash. A checkpoint must be copied to a trust domain the database owner cannot rewrite, such as WORM object storage, a transparency service, or an RFC~3161 timestamp authority~\cite{adams2001timestamp}. A checkpoint left beside the database mainly detects accidents.

Hash chaining detects editing, insertion, reordering, and middle deletion when the verifier has a trusted expected head. It cannot by itself detect deleting a suffix or replacing the entire database and recomputing every hash. The external head closes that gap relative to its anchoring time.

\subsection{Receipts}

A deletion receipt binds subject, time, selector digest, purged record/episode IDs, audit event ID/hash, and receipt digest. An action receipt binds transaction, subject, exact plan hash, terminal state, operation-result/compensation digests, audit event ID/hash, and receipt digest. Verification requires:

\begin{enumerate}
  \item the receipt digest recomputes;
  \item its event exists on a valid subject chain;
  \item the event hash equals the receipt's bound hash;
  \item event payload and receipt agree on transaction, plan, terminal state, and operation digest; and
  \item plan and recorded operational tables recompute consistently.
\end{enumerate}

\begin{figure*}[t]
\centering
\resizebox{0.98\textwidth}{!}{%
\begin{tikzpicture}[
  font=\sffamily\small,
  event/.style={draw=aetnaBlue, rounded corners=2pt, minimum width=27mm, minimum height=13mm, align=center, fill=aetnaBlue!7},
  artifact/.style={draw=aetnaPurple, rounded corners=2pt, minimum width=31mm, minimum height=12mm, align=center, fill=aetnaPurple!8},
  verifier/.style={draw=aetnaTeal, rounded corners=2pt, minimum width=31mm, minimum height=12mm, align=center, fill=aetnaTeal!9},
  arrow/.style={-{Latex[length=2.2mm]}, line width=.8pt, draw=aetnaInk},
  bind/.style={-{Latex[length=2.2mm]}, dashed, line width=.8pt, draw=aetnaPurple},
  tinylabel/.style={font=\sffamily\scriptsize, text=aetnaInk!80, align=center}
]
  \node[event] (e1) {$e_1$\\episode ingested\\$h_1=H(e_1,\bot)$};
  \node[event, right=8mm of e1] (e4) {$e_4$\\record quarantined\\$h_4=H(e_4,h_3)$};
  \node[event, right=8mm of e4] (e6) {$e_6$\\action proposed\\$h_6=H(e_6,h_5)$};
  \node[event, right=8mm of e6] (e7) {$e_7$\\plan approved\\$h_7=H(e_7,h_6)$};
  \node[event, right=8mm of e7] (e12) {$e_{12}$\\receipt issued\\$h_{12}=H(e_{12},h_{11})$};

  \draw[arrow] (e1) -- node[tinylabel, above]{previous hash} (e4);
  \draw[arrow] (e4) -- node[tinylabel, above]{previous hash} (e6);
  \draw[arrow] (e6) -- node[tinylabel, above]{previous hash} (e7);
  \draw[arrow] (e7) -- node[tinylabel, above]{previous hash} (e12);

  \node[artifact, below=17mm of e6] (plan) {WorldPatch\\plan hash $p$};
  \node[artifact, below=17mm of e7] (approval) {Approval\\$\mathrm{HMAC}_k(p,\mathrm{scope},t,n)$};
  \node[artifact, below=17mm of e12] (receipt) {Action receipt\\$H(p,\mathrm{terminal},\mathrm{ops},h_{12})$};
  \node[artifact, right=12mm of e12] (checkpoint) {Checkpoint\\subject head $(12,h_{12})$};
  \node[verifier, below=17mm of checkpoint] (verify) {Standalone verifier\\recompute chain, plan,\\approval, and receipt};

  \draw[bind] (plan) -- (e6);
  \draw[bind] (plan) -- (approval);
  \draw[bind] (approval) -- (e7);
  \draw[bind] (plan) -- (receipt);
  \draw[bind] (receipt) -- (e12);
  \draw[bind] (e12) -- (checkpoint);
  \draw[arrow] (checkpoint) -- (verify);
  \draw[arrow] (receipt) -- (verify);
  \draw[arrow] (plan) |- (verify);

  \node[tinylabel, below=3mm of e4] {Changing \texttt{quarantined} to \texttt{active} changes $h_4$\\and invalidates the recorded suffix/checkpoint.};
\end{tikzpicture}}
\caption{Evidence binding in the recorded flagship run (intermediate events are visually elided, not omitted from verification). Hash chaining detects prefix mutation; an externally held checkpoint is required to detect whole-history replacement or suffix truncation. The construction is tamper-evident, not tamper-proof.}
\label{fig:audit-binding}
\end{figure*}


\subsection{Independent implementations}

The repository ships two Python-standard-library programs. \code{verify\_audit.py} implements the event-chain, checkpoint, and deletion-receipt rules without importing \system. \code{verify\_actions.py} additionally reconstructs \worldpatch, validates approval scope and optional HMAC, and checks action receipt bindings. This is implementation independence: a defect in the engine's self-check need not be repeated by calling the same library. It is not organizational independence; both programs are maintained in the same project.

\property{E1 (Mutation detection relative to an anchor)} If an adversary modifies an event at or before a trusted checkpoint's pinned head without finding a hash collision, then either the chain rule or checkpoint containment check fails.

\emph{Argument.} Modifying event $e_j$ changes $h_j$ by \cref{eq:eventhash}. Retaining the old suffix breaks the next \code{prev\_hash}; recomputing the suffix changes the pinned head. Deleting the suffix fails containment. Replacing the external checkpoint lies outside the assumption.

\subsection{Evidence is not truth}

The ledger proves consistency among recorded assertions. It does not make a false adapter observation true, authenticate a forged host source label, or prove that no unlogged side channel existed. In particular:

\begin{itemize}
  \item \emph{integrity} means recorded bytes have a consistent hash history;
  \item \emph{authenticity} depends on keys, host identity, and adapter/provider evidence;
  \item \emph{completeness} depends on mediation and event coverage; and
  \item \emph{semantic correctness} depends on policy and verifiers.
\end{itemize}

This distinction avoids treating ``blockchain-like'' logging as a substitute for access control or trustworthy observation.
